\chapter{Derivazione delle distribuzioni binomiale negativa e logaritmica}
Nel modello parametrico analizzato in questo lavoro, assumeremo che la SAD sia ben descritta da una binomiale negativa.
Prima di entrare nel merito dei modelli per la stima del numero di specie, vediamo come si può ricavare tale distribuzione (come anche la distribuzione logaritmica, molto utilizzata in ecologia), modellando la dinamica stocastica dell'abbondanza delle specie attraverso la così detta \emph{birth-death master equation}\cite{2016AzaeleSuweis}.\\
Entrambe le distribuzioni (logaritmica e binomiale negativa) possono essere derivate da processi ecologici fondamentali quale appunti nascita, morte e migrazione. 

Sia $\emph{P}_{\emph{n,s}}(t)$ la probabilità che ad un certo tempo $t$, la specie \emph{s} abbia esattamente \emph{n} individui, con $\emph{s}\in\left \{ 1,...,S \right \}$. Assumiamo che la dinamica della popolazione di ogni specie $s=1,2,...,S$ con popolazione $n$ sia governata dai rates di nascita e di morte, rispettivamente, $\emph{b}_\emph{{n,s}}$ e $\emph{d}_\emph{{n,s}}$.
L'equazione che regola l'evoluzione di $\emph{P}_{\emph{n,s}}(t)$ per $\emph{n}\ge0$ sarà quindi:
\begin{equation}
\frac{\partial\emph{P}_{\emph{n,s}}(t)}{\partial\emph{t}}=
\emph{P}_{\emph{n-1,s}}(t)\emph{b}_\emph{{n-1,s}}+\emph{P}_{\emph{n+1,s}}(t)\emph{d}_\emph{{n+1,s}}-\emph{P}_{\emph{n,s}}(t)\emph{b}_\emph{{n,s}}-\emph{P}_{\emph{n,s}}(t)\emph{d}_\emph{{n,s}}.
\label{eq:master}
\end{equation}
In questo lavoro, per evitare che \emph{n} sia negativo, imponiamo delle condizioni al contorno riflettenti: $\emph{b}_\emph{{-1,s}}=\emph{d}_\emph{{0,s}}=0$. %la (\ref{eq:master}) è valida anche per $\emph{n}=0$ e \emph{n=1}. 
Si trova che per $\emph{n}>0$ la soluzione stazionaria è:
\begin{equation}
\emph{P}_{\emph{n,s}}=P_{0,\emph{s}}\prod_{i=0}^{n-1}\frac{\emph{b}_\emph{{i,s}}}{\emph{d}_\emph{{i+1,s}}}
\label{eq:steadystate}
\end{equation}
dove il termine $P_{0,\emph{s}}$ è il fattore di normalizzazione che può essere trovato imponendo la condizione $\sum_{n=1}^\infty \emph{P}_{\emph{n,s}}=1.$ Notiamo che la somma inizia da \emph{n}=1 in quanto non si considerano specie con abbondanza nulla.\\



\section{Distribuzione binomiale negativa}
Ipotizziamo che tutti gli individui, siano essi appartenenti a specie rare o comuni, abbiano la stessa probabilità di morire, sopravvivere e riprodursi. In questo caso i tassi di nascita e morte pro capite non dipendono dal numero di individui \emph{n} appartenenti alla specie.
Definiamo quindi $\emph{b}_\emph{{n,s}}$ come:
\begin{equation}
\emph{b}_\emph{n,s}=\emph{b}_\emph{s}(n+\emph{r}_\emph{s}),
\label{eq:birthNB}
\end{equation}
dove $\emph{r}_\emph{s}$ è un parametro che tiene conto di eventi di immigrazione o di interazioni intraspecifiche.\newline
Analogamente definiamo $\emph{d}_\emph{n,s}$: 
\begin{equation}
\emph{d}_\emph{n,s}=\emph{d}_\emph{s}\emph{n}.
\label{eq:deathNB}
\end{equation}

Sostituendo questi ultimi termini nella (\ref{eq:steadystate}) e denotando con $\xi_\emph{s}=\emph{b}_\emph{s}/\emph{d}_\emph{s}$, si ottiene:
$$
\emph{P}_\emph{n,s}=P_\emph{0,s}\binom{n+\emph{r}_\emph{s}-1}{n}\xi_\emph{s}^n.
$$
La costante di normalizzazione può essere calcolata imponendo:
$$
1=\sum_{n=1}^\infty \emph{P}_{\emph{n,s}}=P_\emph{0,s}\sum_{n=1}^\infty\binom{n+\emph{r}_\emph{s}-1}{n}\xi_\emph{s}^n=P_\emph{0,s}[1-(1-\xi_\emph{s})^{\emph{r}_\emph{s}}](1-\xi_\emph{s})^{-\emph{r}_\emph{s}}.
$$
Dunque la probabilità che una specie \emph{s} abbia \emph{n} individui all'equilibrio è data da una binomiale negativa di parametri $(\emph{r}_\emph{s}, \xi_\emph{s})$ e normalizzata per abbondanze non nulle ($n\ge 1$):
\begin{equation}
\emph{P}_\emph{n,s}^{\emph{NB}}=\frac{1}{1-(1-\xi_\emph{s})^{\emph{r}_\emph{s}}}\binom{n+\emph{r}_\emph{s}-1}{n}\xi_\emph{s}^n(1-\xi_\emph{s})^{\emph{r}_\emph{s}}.
\label{eq:NBprobability}
\end{equation}
Sotto l'ipotesi della teoria neutrale secondo la quale le specie sono considerate demograficamente equivalenti, possiamo rimuovere l'indice \emph{s} di specie dall'equazione sopra, ottenendo così una SAD per l'ecosistema in esame. Notiamo che quindi in questo framework, la dinamica della popolazione della specie è determinata puramente da processi demografici random (nascita, morte e migrazione) e ogni specie rappresenta una diversa realizzazione dello stesso processo stocastico.

\section{Distribuzione logaritmica di Fisher}
Notiamo che, scegliendo in modo diverso il termine $\emph{b}_\emph{n,s}$, si può ottenere semore attraverso la \emph{birth-death master equation} (\ref{eq:master}), un'altra importante distribuzione: la distribuzione logaritmica di Fisher.
Consideriamo ora nella dinamica della comunità ecologica anche il fenomeno di speciazione (cioè l'entrata di nuove specie nel sistema dovute a mutazioni, invece che a migrazione da comunità esterne). La speciazione avverrà con un rate molto piccolo $\nu$, e si avrà:
\begin{equation}
    \emph{b}_\emph{n,s}=\emph{b}_\emph{s}n+\delta_\emph{n,0}\nu.
\label{eq:birthlog}
\end{equation}
Aggiungendo la condizione al contorno riflettente $\emph{b}_\emph{0,s}=\nu$ si ha quindi che il tasso di nascita tiene conto della riproduzione e della speciazione. In particolare, il parametro $\nu$ assicura che, se le specie si estinguono, la comunità rimane sempre popolata da un individuo.
Dunque sostituendo la (\ref{eq:deathNB}) e la (\ref{eq:birthlog}) nella (\ref{eq:steadystate}) e definendo $\emph{x}_\emph{s}=\emph{b}_\emph{s}/\emph{d}_\emph{s}$, si trova la seguente soluzione stazionaria:
\begin{equation}
    \emph{P}_\emph{n,s}=P_\emph{0,s}\frac{\nu}{\emph{b}_\emph{s}}\frac{\emph{x}_\emph{s}^{\emph{n}}}{\emph{n}}.
\end{equation}
La costante di normalizzazione $P_\emph{0,s}$ si determina imponendo:
$$
1=\sum_{n=1}^\infty \emph{P}_{\emph{n,s}}=P_\emph{0,s}\frac{\nu}{\emph{b}_\emph{s}}\sum_{n=0}^\infty\frac{\emph{x}_\emph{s}^{\emph{n}}}{\emph{n}}=P_\emph{0,s}\frac{\nu}{\emph{b}_\emph{s}}[-\log(1-\emph{x}_\emph{s})].
$$
Dunque abbiamo:
\begin{equation}
    \emph{P}_\emph{n,s}^{\emph{LS}}=-\frac{1}{\log(1-\emph{x}_\emph{s})}\frac{\emph{x}_\emph{s}^{\emph{n}}}{\emph{n}}.
\label{eq:fisherdist}
\end{equation}
Anche in questo caso assumiamo che le specie siano equivalenti e possiamo dunque omettere l'indice \emph{s}.

\subsection{La distribuzione di Fisher come caso particolare della binomiale negativa}
Osserviamo che la distribuzione binomiale negativa converge ad una distribuzione logaritmica nel limite di \emph{r} che tende a zero:
\begin{equation}
    \lim_{\emph{r} \to 0}\emph{P}_\emph{n}^{\emph{NB}}= \lim_{\emph{r}\to0}\frac{(1-\xi)^{\emph{r}}}{1-(1-\xi)^{\emph{r}}}\binom{n+\emph{r}-1}{n}\xi^n=\frac{\xi^n}{-n\ln(1-\xi)},
\label{eq:convergence}
\end{equation}
dove si è usato il fatto che:
$$
\binom{n+\emph{r}-1}{n}=\frac{\Gamma(\emph{n}+\emph{r})}{\Gamma(\emph{n+1})\Gamma(\emph{r})} \approx\frac{\emph{r}}{\emph{n}},
$$
per $\emph{r}\approx 0$.\\
Notiamo dunque che la (\ref{eq:convergence}) coincide con la (\ref{eq:fisherdist}) ponendo  $\emph{x}=\xi$.


%\section{Protocolli crittografici}
%HTTPS, TLS, Diffie-Hellman, SHA-1
% ---- ELEMENTI UTILI E GIA' PRONTI! ----
%Secondo capitolo della tesi. Esempio di citazione doppia \cite{Munoz-Lipo,Vas}.

%Esempio di figura in \figurename\ \ref{FIG:LogoUniPD}.
%
%\begin{figure}[!htbp]
%\centering
%\includegraphics[width=0.25\textwidth]{./figure//LogoUniPD}
%\caption{Esempio di figura.}
%\label{FIG:LogoUniPD}
%\end{figure}
%
%Esempio di tabella in \tablename\ \ref{TAB:Esempio}.
%
%\begin{table}[!htbp]
%\centering
%\renewcommand{\arraystretch}{1.3}
%\caption{Esempio di tabella.}
%\begin{tabular}{cc}
%\hline
%Nome & Valore \\
%\hline
%a & 1 \\
%b & 2 \\
%c & 3 \\
%d & 4 \\
%e & 5 \\
%f & 6 \\
%\hline
%\end{tabular}
%\label{TAB:Esempio}
%\end{table}